% chapters/ch14.tex —— 第 14 章
% 体例：具体失败 → 设计问题 → 机制解释 → 对照实验 → 条件变化 → 选择原则
% 数字身份纪律：
%   - E02 运行 E02-20260924T153511Z-421eac 的噪声臂与 E01 的数字一律 \measured{}；
%   - 2607.07663（RSI survey §5.3）、2603.11768（SSGM）的数字与定义一律 \cited{}；
%   - 由上述数字算出的倍数与差值用 \calcnum{}。
% P6 依赖 E12（未跑），本章一律写"待补"，不编造退化量。
\chapter[可靠性、隔离与可恢复的更新]{可靠性、隔离与\\可恢复的更新}
\label{ch14}

\chaptergoal{一次有害的持久修改，怎样被及时发现，又怎样被撤回？}
\chapterprops{P6关注哪些更新引起旧能力退化，以及怎样保持和恢复。退化量统一为$D=\text{旧版本得分}-\text{新版本得分}$，$D>0$表示退化；E12尚无真实模型结果。P8作为反馈接地频率的扩展问题，依赖尚未执行的E02b。}
\chapterfig{图~\ref{fig:ch14-rail}：可变轨道 / 不可变轨道与隔离边界}
\chapterpitfall{同一个 Agent 修改了实现、测试与评分逻辑，然后得到一个漂亮的提升——
每一步单独看都合法}

% ---------- 1. 具体失败 ----------
\begin{failurebox}
E02在训练阶段设置了一个噪声臂，用来观察错误反馈下的表现：写记忆时所见的外部判定，
有 \measured{2} 条被我们人为翻转（\measured{8} 条反馈里翻了 \measured{2} 条，
翻转率 \measured{0.2500}）。提示未显式标明翻转位置；这不证明系统没有察觉标签与执行输出的矛盾。

这一臂的封存题回合通过率是 \measured{0.6625}，相对无记忆基线 \measured{0.8125}，
收益 \measured{$-0.1500$}，95\% 区间 \measured{$[-0.3125, +0.0125]$}。
在 \measured{16} 道配对题里，它比基线差的题有 \measured{8} 道、好的只有 \measured{3} 道。

\textbf{标签格式本身不能识别真假。}旧流程将判定前缀翻转后仍保留执行输出，见第\ref{ch04}章的轨迹核查；错误标签与其他证据可能不一致。是否察觉、写成什么经验以及后来是否取回，需要逐层核对，不能从“翻转了两条”直接推断“成功植入了两条错误记忆”。该流程没有提供一个经独立校准的污染报警率。
\end{failurebox}

\designquestion{前面各章处理的是“怎么让系统变好”，本章处理的是另一半：
\textbf{当一次更新是有害的，靠什么发现它、靠什么把它退回去、以及怎样降低同类问题再次发生的风险}。
这三个问题各自需要不同的机制：发现要可观测，退回要可恢复，不再犯要可归属。}

第 2 章把 $\mathcal V$ 写成公理时留了一个工程后果没有展开：
候选若能任意修改验收规则，前后分数可能不再可比；候选若能读取最终测试内容，则涉及评价信息泄漏。可改与可读是两种不同边界，公开的开发测试可以合理用于排错，却不能因此被称为隐藏评价。
本章从污染来源与权限边界出发，说明双轨记录、回归评价与恢复流程，再讨论长期漂移、闭环退化及其检验设计。来源链用于定位问题，恢复记录用于确认撤回了什么，独立行为评价用于判断旧能力是否改变。

% ---------- 2. 污染入口 ----------
\section{不可信反馈如何进入持久状态}
\label{sec:ch14-poison}

\begin{sloppypar}
本节先讨论三个入口：反馈本身错误、反馈被篡改、正确反馈在写入时失真。它们不是穷尽分类，过期信息、适用范围丢失和检索串用也可能带来错误。三种入口都可能留下格式正常的记录，具体症状仍需沿来源链识别。
\end{sloppypar}

\begin{table}[htbp]
\centering
\caption{三种污染入口：症状、自查可能受限的原因、本书的防线}
\label{tab:ch14-poison}
\begin{tabular}{>{\raggedright\arraybackslash}p{2.6cm}>{\raggedright\arraybackslash}p{3.8cm}>{\raggedright\arraybackslash}p{4.0cm}>{\raggedright\arraybackslash}p{3.0cm}}
\toprule
入口 & 症状 & 自查可能受限的原因 & 本书的防线 \\
\midrule
反馈本身错 & 置信度高的错误被反复强化 & 共享模型可能有相同盲点，相关性需检验 & 反馈源分级；接地优先 \\
\midrule
反馈被篡改 & 翻转后的判定与正常判定同形 & 仅看标签格式无法区分；执行线索仍可比对 & 原始记录受保护；保留来源 \\
\midrule
写入时失真 & 摘要把事实改写掉了 & 若未保存原文，事后难以核对压缩失真 & 双轨：原文只追加，摘要可重算 \\
\bottomrule
\end{tabular}
\end{table}

\begin{sloppypar}
RSI 综述将共享生成与评估机制可能形成的错误强化归入“自证循环”，并回顾自奖励强化学习等具体研究\cite{rsi-survey-2607}。这是综述的归纳，不是本书直接复现的因果效应。共享权重提供了共同盲点的一种可能来源，但不能据此断言任意同模型评审都必然相关、都会优先奖励高置信度错误。检查评委与生成器的错误重合、独立依据和校准，比仅把提示词改得更严厉更有诊断价值。换模型也不自动带来独立证据，因为任务材料、训练来源和评价目标仍可能相同。
\end{sloppypar}

\begin{sloppypar}
E02为第二类入口保留了一次真实运行记录：噪声臂把 \measured{2}/\measured{8} 条外部判定翻转，
得到的收益点估计 \measured{$-0.1500$}、区间 \measured{$[-0.3125, +0.0125]$}；
把噪声臂与未翻转判定的测试反馈臂相比，点估计 \measured{$-0.0500$}、
区间 \measured{$[-0.2000, +0.1125]$}。
两个区间都跨 0（可用配对题 \measured{16} 道），
所以\textbf{点估计偏负，但现有区间尚未确定所测方案差异的方向}。
各臂分别完成训练与经验构建，来源验收结果及最终记忆并不相同。这组比较包含反馈、候选与记忆路径的共同变化，不能把两臂点差单独归因于两条标签翻转。若要估计更窄的标签效应，应固定同一组来源输出，再按预定规则改变反馈；这也仍需区分错误标签被呈现、被写入和被后续任务采用。
\end{sloppypar}

% ---------- 3. 隔离 ----------
\subsection{从一条错误记录追到受影响的任务}
\label{sec:ch14-provenance-chain}

持久污染的难点在于它跨越多个时点。错误出现时未必影响当前任务，后来被取回时又可能已经看不出原来依据。因此一次可靠性调查不能只问“这条记忆对不对”，还应问它从哪里来、经过哪些改写、在哪些版本中发布、被哪些任务实际读取。本节给出本书用于组织调查的来源链，不宣称它已经消除了错误。

\subsubsection{原始证据、解释与建议分别保存}

一份反馈至少可能包含三类内容：工具执行的原始输出、对输出原因的解释、建议未来采取的动作。执行返回非零状态是一个观察；“失败因为边界少算一次”是解释；“以后都在上界加一”是建议。三类内容的可靠性不同，不能因为它们一起出现在一段文本里，就继承同一个可信标签。

E02 的标签扰动说明这种拆分的重要性。判定前缀被翻转，执行输出仍可能提供相反线索。模型可能信前缀，也可能信输出，或者写出对两者都不敏感的通用建议。只检查提示里存在错误标签，可以确认输入受到扰动；检查最终经验，才知道错误内容是否发布；检查后续任务的检索载荷与行为，才知道它是否进入决策。这三项证据各自需要入口，不能在汇总表中用一个“污染成功”覆盖。

来源记录可以保存任务标识、执行器版本、反馈产生时刻、原始输出摘要值与完整原文位置。派生条目另外保存构建提示、候选文本、来源引用和准入状态。若为了保护敏感内容而不保留全文，也要明确哪些重建能力因此丢失，不能在缺原文的情况下承诺逐字追溯。本书当前实验使用合成任务，不据此推定任意真实用户数据都应无限期保存。

解释经过多轮压缩后，最容易丢失的是限定条件。例如来源只观察到“空列表时返回默认值”，第一次摘要变成“缺值时用默认值”，第二次又变成“任何错误都返回默认值”。每一步看起来都在概括，合起来却改变了异常契约。检查最终句子是否流畅没有帮助，应比较关键条件、对象、否定词与例外是否仍能在来源中找到支持。

如果一个条目引用多个来源，也不宜把它们压成一个无方向的列表。哪些来源支持事实、哪些只是类似任务、哪些提供反例，可以分别标记。发现某个来源错误时，系统才能定位受影响的派生条目。来源依赖图可以帮助扩大检查范围，但“依赖受污染来源”只说明需要复查，不自动证明所有下游结论都错；有些结论可能还有独立证据支持。

\subsubsection{有来源的内容仍可能不适用}

追溯解决的是内容从何而来，不是它在当前情形是否适用。某条经验忠实复述了过去版本的接口，却可能在接口升级后失效；某个用户的偏好记录真实，也可能不属于当前用户；某次环境故障确实发生过，却不适合作为所有后续任务的永久规则。这些都不需要源材料被篡改就会造成错误。

因此读取时需要核对作用域。任务族、工具版本、用户或项目身份、有效时间与依赖条件，可以作为候选检索之后的约束。语义相似度高只表示表述相关，不表示访问被允许、条件相同或信息仍有效。SSGM 的读取公式将语义候选、访问规则与时效约束分开\cite{ssgm-2603}；这种分工有助于审查，但访问规则与时效函数本身仍需实现和测试。

时效也不等于简单的“越新越可信”。一条新的误报可能不如一条旧的、经过多次独立验证的规则可靠。可以区分事实何时发生、何时被观察、何时被写入与何时失效；系统在重写时不应仅更新存储时间，就把一个旧事实伪装成新证据。发生冲突时，先保留版本关系和适用条件，通常比自动用最新句子覆盖旧句子更便于诊断。

E12 v2 的快照组件把文本和来源元数据分别比较，也是为了避免混淆。“文字未变但来源换了”可能改变可信解释；“文字改了但字节数没变”仍可能改变行为。组件能报告这些事实，却不根据变更大小自动判好坏。最终影响还需回到固定评价任务与实际读取记录。

\section{修改空间与评测环境的隔离}
\label{sec:ch14-iso}

\begin{sloppypar}
公理 1 说 $\mathcal V$ 默认冻结。它的工程表达是一条\textbf{单向的读写边界}（图~\ref{fig:ch14-rail}）：
\end{sloppypar}

% figures/ch14-rail.tex —— 第 14 章核心图（全书引用入口）
% 用法：在 book/chapters/ch14.tex 中 % figures/ch14-rail.tex —— 第 14 章核心图（全书引用入口）
% 用法：在 book/chapters/ch14.tex 中 % figures/ch14-rail.tex —— 第 14 章核心图（全书引用入口）
% 用法：在 book/chapters/ch14.tex 中 \input{figures/ch14-rail}
% 图主体在 ch14-rail-body.tex，便于单独编译与复用。
\begin{figure}[htbp]
\centering
\input{figures/ch14-rail-body}
\caption{来源、可变状态、开发反馈与独立评价的职责分工。保留原始记录支持追溯，历史快照支持指定状态恢复；测试可读与规则可改是两种不同权限。图示是本书的工程设计建议，不声称旧实验已实现全部隔离。}
\label{fig:ch14-rail}
\end{figure}

% 图主体在 ch14-rail-body.tex，便于单独编译与复用。
\begin{figure}[htbp]
\centering
% Proposed responsibility boundaries; not a claim of deployed OS isolation.
\begin{tikzpicture}[x=1cm,y=1cm,
  block/.style={draw,rounded corners=2pt,text width=3cm,minimum height=1cm,align=center,font=\small},
  arrow/.style={-{Stealth[length=1.7mm]},thick},
  note/.style={text width=13cm,align=left,font=\footnotesize}]
  \node[block,fill=blue!8] (raw) at (0,2) {只追加原始记录\\来源、反馈、历史版本};
  \node[block,fill=orange!10] (memory) at (0,-0.2) {可变状态与快照\\摘要、技能、检索索引};
  \node[block,fill=orange!10] (candidate) at (5,2) {候选版本\\绑定来源与依赖清单};
  \node[block,fill=gray!8] (dev) at (5,-0.2) {开发测试与反馈\\允许排错的具体范围};
  \node[block,fill=green!9] (judge) at (10,2) {独立评价流程\\固定规则与版本记录};
  \node[block,fill=green!9] (result) at (10,-0.2) {结果与发布决定\\报告实际信息边界};
  \draw[arrow] (raw) -- (candidate);
  \draw[arrow] (memory.east) -- (candidate.south west);
  \draw[<->,>=Stealth,thick] (candidate.south) -- (dev.north);
  \draw[arrow] (candidate) -- node[above,font=\scriptsize] {提交工件} (judge);
  \draw[arrow] (judge) -- (result);
  \draw[dashed,thick] (7.5,-1.1) -- (7.5,3.1);
  \node[font=\scriptsize] at (10,3.05) {评价规则不由候选直接改写};
  \draw[arrow] (raw.south) -- node[right,font=\scriptsize] {追溯、重建} (memory.north);
  \node[note,anchor=north west] at (-1.5,-1.6) {图示为设计分工：开发反馈可以可读；最终评价若声明隐藏，须另有实际访问隔离。\\本书旧实验的测试写入限制与查询账不构成完整隐藏测试，也不构成操作系统级沙箱。};
\end{tikzpicture}

\caption{来源、可变状态、开发反馈与独立评价的职责分工。保留原始记录支持追溯，历史快照支持指定状态恢复；测试可读与规则可改是两种不同权限。图示是本书的工程设计建议，不声称旧实验已实现全部隔离。}
\label{fig:ch14-rail}
\end{figure}

% 图主体在 ch14-rail-body.tex，便于单独编译与复用。
\begin{figure}[htbp]
\centering
% Proposed responsibility boundaries; not a claim of deployed OS isolation.
\begin{tikzpicture}[x=1cm,y=1cm,
  block/.style={draw,rounded corners=2pt,text width=3cm,minimum height=1cm,align=center,font=\small},
  arrow/.style={-{Stealth[length=1.7mm]},thick},
  note/.style={text width=13cm,align=left,font=\footnotesize}]
  \node[block,fill=blue!8] (raw) at (0,2) {只追加原始记录\\来源、反馈、历史版本};
  \node[block,fill=orange!10] (memory) at (0,-0.2) {可变状态与快照\\摘要、技能、检索索引};
  \node[block,fill=orange!10] (candidate) at (5,2) {候选版本\\绑定来源与依赖清单};
  \node[block,fill=gray!8] (dev) at (5,-0.2) {开发测试与反馈\\允许排错的具体范围};
  \node[block,fill=green!9] (judge) at (10,2) {独立评价流程\\固定规则与版本记录};
  \node[block,fill=green!9] (result) at (10,-0.2) {结果与发布决定\\报告实际信息边界};
  \draw[arrow] (raw) -- (candidate);
  \draw[arrow] (memory.east) -- (candidate.south west);
  \draw[<->,>=Stealth,thick] (candidate.south) -- (dev.north);
  \draw[arrow] (candidate) -- node[above,font=\scriptsize] {提交工件} (judge);
  \draw[arrow] (judge) -- (result);
  \draw[dashed,thick] (7.5,-1.1) -- (7.5,3.1);
  \node[font=\scriptsize] at (10,3.05) {评价规则不由候选直接改写};
  \draw[arrow] (raw.south) -- node[right,font=\scriptsize] {追溯、重建} (memory.north);
  \node[note,anchor=north west] at (-1.5,-1.6) {图示为设计分工：开发反馈可以可读；最终评价若声明隐藏，须另有实际访问隔离。\\本书旧实验的测试写入限制与查询账不构成完整隐藏测试，也不构成操作系统级沙箱。};
\end{tikzpicture}

\caption{来源、可变状态、开发反馈与独立评价的职责分工。保留原始记录支持追溯，历史快照支持指定状态恢复；测试可读与规则可改是两种不同权限。图示是本书的工程设计建议，不声称旧实验已实现全部隔离。}
\label{fig:ch14-rail}
\end{figure}


\begin{itemize}
  \item \textbf{候选不能改}。任务划分、验收器、最终测试集三者都不在 Updatable 内。
        把 $\mathcal V$ 纳入可更新集合需要单列实验组，并且必须回答“最终报告从哪来”。
  \item \textbf{读取边界须单独实现}。最终评价材料应与开发反馈区分。本书旧实验台的 protected 规则主要限制工具写入，原始 E02 轨迹中确有读取测试文件的操作。\texttt{QueryBudget} 记录 sealed 查询并提供事后检查，不是文件访问拦截器；调用者也必须实际执行检查。不能把这些组件合称为已经实现隐藏测试。
  \item \textbf{执行器的权限要与声明匹配}。任务目录和工具白名单能减少误操作，但执行生成的 Python 代码仍不等于操作系统级隔离。静态 Agent 的 \texttt{persistent\_writes} 断言只检查被该字段登记的持久写入，不证明文件系统与网络没有其他副作用。
\end{itemize}

\begin{pitfallbox}
\textbf{反例：把验收器一起改掉之后报告提升。}
一个 Agent 若被授予实现、测试与评分逻辑的写权限，它有三条能改变分数的路径
（改实现 / 改测试 / 改评分）都能让数字变好。
本书要求把最终验收规则放在候选的可修改范围之外，并把实际改动写入记录；现有实验台的工具级限制不能替代更强的权限隔离。
区分不了这三者的提升，一律不算收益。
\end{pitfallbox}

% ---------- 4. 双轨 ----------
\subsection{隔离需要覆盖哪些具体权限}
\label{sec:ch14-permission-boundary}

“在沙箱里运行”是范围很宽的说法。它可能只表示给任务一个临时目录，也可能表示进程无法读取其他目录、不能访问网络、只能执行受限命令。不同实现的约束强度差别很大。书中应描述已经实施的边界，而不能用同一个词暗示完整安全保证。

\subsubsection{工具接口限制与进程限制}

本书旧文件工具会检查部分写入路径和测试文件名，这是通过指定工具修改文件时的限制。如果候选代码之后作为普通 Python 程序执行，代码可能调用标准库读写文件或启动子进程；先前工具函数上的检查未必覆盖这些路径。因此“工具拒绝写测试”与“所有执行都不可能影响测试”是两个不同主张。前者可由工具检查验证，后者需要更强的进程权限与资源隔离。

新版 E11 为自身工具路径增加了任务目录归属检查，并把可写文件限于初始可变源码，保留旧实现不改。这降低了该工具接口的误写范围，但正文仍明确它不是操作系统级安全隔离。实验使用受控的合成编程任务，并不能把这一实现推广成处理任意不可信程序的生产安全方案。

进一步的本地检查给出了具体反例：旧通用文件组件把“字符串以前缀开头”误当成“路径是根目录后代”，相邻同名前缀目录因而可能通过检查。受控子进程也能在绕开该文件工具后修改一个虚构的受保护测试标记。复现只使用新建临时文件，没有触碰历史实验或真实隐藏材料；其价值在于证伪全范围保护的说法。E01独立新版文件边界已经修正清单与路径检查，但候选代码的进程权限仍须另行实现，不能将这次局部修复写成安全隔离全部完成。

读取限制同样需要落实到实际入口。最终测试文件若已经复制到候选可读目录，外层的查询计数对象无法阻止候选读它；即使没有显式读取，测试错误输出也可能逐步暴露断言。使用开发测试排错并非错误，关键是如实定义评价条件，并把真正留出的测试与开发反馈分开。隐藏测试的声明需要同时检查文件暴露、工具返回、日志访问和测试服务接口。

评价器自身也有依赖。测试命令、运行时版本、依赖包和评分逻辑的变化，都可能改变结果。候选文件未动而分数变高，可能来自环境升级；测试文件未动而行为变了，也可能来自导入路径或依赖替换。将这些对象写入运行清单，比只保存候选源码更接近可复查的比较。本书新版实验的源与任务摘要值冻结，主要服务于这种一致性核对，并不证明全部环境都可复现。

\subsubsection{评估反馈与最终判断的权限}

一个智能体完全可以执行开发测试、读取报错、提出测试补充，也可以参与生成候选验收用例。问题在于谁批准评分规则进入最终比较，以及规则变化后如何保持参照。若智能体为了某个候选修改测试，应把这视为验证器共同演化的一部分，另留独立检查；不能把同一次测试修改隐藏在“实现优化”里。

更细的权限划分可以允许提出修改而禁止直接生效。候选将新测试写入待审工件，独立流程检查后形成新的验证器版本，再决定是否用于下一轮实验。旧版结果仍按旧版规则解释，不能用新版评分覆盖旧报告而不保留映射。这里的独立可以指职责、输入或执行权限上的分离；同一模型换一个角色提示，并不自动等于独立证据。

还要区分查询预算与访问授权。预算回答某个允许的评价最多调用多少次；授权回答它是否可以被当前主体调用。把次数设为零但仍允许读取本地测试文件，无法实现隐藏；把访问隔离做好却无限重复看同一测试反馈，也可能导致选择过拟合。两种控制需要同时检查，任何一个单独存在都不支持全部评价纪律已经落实。

\section{双轨设计：可变轨道与不可变轨道}
\label{sec:ch14-rail}

\begin{sloppypar}
让基于来源的重建成为可能的先决条件是：\textbf{原始记录还在}。恢复指定旧版本则还可依靠完整历史快照。
一份被反复摘要重写过的记忆，即使知道它错了，也无法恢复出它当初依据的是什么。
双轨设计因此把持久物分成两条轨道\cite{ssgm-2603}：
\end{sloppypar}

\begin{itemize}
  \item \textbf{不可变轨道}：只追加的原始情景记录（原文记作 $K_{\text{ledger}}$）。
        在该设计的常规读取公式中，推理从可变状态取上下文，原始账用于校正与追溯；也可另行设计受控的原文检索。
  \item \textbf{可变轨道}：结构化记忆、摘要、技能库。它服务于推理，可以被改写、被淘汰。
  \item \textbf{两道门}：写前的\emph{写校验门}（先与既有核心事实做矛盾检查再提交），
        读时的\emph{读过滤门}（按作用域与时效放行）。
  \item \textbf{异步对账}：可变轨道定期与不可变轨道重算对齐，
        因此摘要被写坏之后可以重算而不是修补。
\end{itemize}

\begin{sloppypar}
这是一份治理框架与研究议程，本书采用其中的分工来组织工程建议，不能称全书实验已经部署双轨系统。可变记忆应能追溯到来源，来源记录的保存与访问应由外部策略约束；回滚需要具体的历史快照或可验证的重建过程。只保留原始记录并不保证可以精确重现某次随机摘要，也不保证原始内容全都正确。SSGM 区分受保护核心事实 $M_{\mathrm{core}}$、操作用原始账 $K_{\mathrm{ledger}}$ 与理想语义目标 $K_{\mathrm{true}}$，三者不能混成一个自动可信的“真值库”\cite{ssgm-2603}。
\end{sloppypar}

% ---------- 5. 回退与放量 ----------
\section{回归、分阶段放量与回退}
\label{sec:ch14-rollback}

\begin{sloppypar}
本书用三个环节组织发布检查，具体采用哪些环节仍依场景与风险而定；它们提供不同的观察机会：
\end{sloppypar}

\begin{table}[htbp]
\centering
\caption{三道发布关：挡什么、判据、落点}
\label{tab:ch14-gate}
\begin{tabular}{>{\raggedright\arraybackslash}p{2.4cm}>{\raggedright\arraybackslash}p{3.8cm}>{\raggedright\arraybackslash}p{4.4cm}>{\raggedright\arraybackslash}p{2.8cm}}
\toprule
关口 & 挡什么 & 判据 & 落点 \\
\midrule
回归集 & 旧能力被换走 & 下降量区间与预定损失容忍线 & \texttt{regression} 划分 \\
\midrule
影子运行 & 新版本在真实分布上出事 & 输出不对外生效，资源仍计账 & \texttt{shadow} \\
\midrule
灰度放量 & 故障影响扩至全量 & $r$ 从小到大，逐级观察 & \texttt{accept($r$)} \\
\bottomrule
\end{tabular}
\end{table}

\begin{sloppypar}
回退这件事本身也有度量。
一份代码 Agent 方向的综述\cite{coding-survey-2608}把\emph{组件退役}与\emph{回滚成功率与成本}列为此领域应当报告的度量，
这些属于它建议补充的评价维度，不能由这段建议推出所有已有方法都未报告相关数值。
本书引用时只取“应当报告”这一层，不补数字；
第 8 章的生命周期成本需要计入维护与恢复，本章 P6 需要对应退化比较；E05 与 E12 各自提供不同证据，不能互相替代（见 \S\ref{sec:ch14-claims}）。
\end{sloppypar}

\subsection{四种``撤回''：程序、状态、效果与补偿}
\label{sec:ch14-undo-kinds}

``回滚''这个词在自进化系统里至少指四件不同的事，
本章章首是反馈标签扰动，并没有执行回退。下面用工程分解说明，为什么恢复文件不等于撤销外部效果，不把它写成已经发生的本书实验事故。

\begin{table}[htbp]
\centering
\caption{四种撤回动作：撤什么、能不能自动、撤完之后剩什么}
\label{tab:ch14-undo}
\begin{tabular}{>{\raggedright\arraybackslash}p{2.4cm}>{\raggedright\arraybackslash}p{3.6cm}>{\raggedright\arraybackslash}p{2.2cm}>{\raggedright\arraybackslash}p{4.8cm}}
\toprule
撤回 & 撤的是什么 & 能否自动 & 撤完之后还剩什么 \\
\midrule
程序回滚 & 持久物本身回到上一个快照 & 依赖完整快照与加载条件 & 被这次更新影响过的\emph{后续决策}不会自动重来 \\
状态回滚 & 进程内上下文回到更新前 & 有检查点时可尝试 & 已经写进持久物的部分不受影响 \\
效果撤销 & 已经作用到外部环境的结果 & 依赖外部接口 & 有些可逆，有些只能补偿；须单列权限和范围 \\
补偿动作 & 用一次新的相反动作抵消 & 部分能 & 留下两笔记录：原动作与补偿动作 \\
\bottomrule
\end{tabular}
\end{table}

\begin{principlebox}
\textbf{程序回滚需要恢复目标工件及其加载条件，字节比对只能确认所检查的文件内容一致。}
版本指针、依赖与运行中状态仍可能需要处理；
但这次更新在它生效期间已经产生的结果——改过的文件、写过的记录、发给外部的动作——
\textbf{不会因为持久物回退而回退}。
表~\ref{tab:ch14-undo} 的最后一列列出了恢复后仍需处理的对象：
判据不是“能不能撤回”，而是“撤回之后还剩什么”。
\end{principlebox}

\begin{sloppypar}
当前 E12 v2 协议第4—5节分别规定工程回退线与恢复评价（\path{experiments/protocols/E12-v2.md}）。文件层检查目标快照的字节与哈希是否一致；行为层另报固定题复评分数及其相对基线、同期对照的差。四个更新臂都安排恢复，分母不只取触发臂。
本计划只测持久物恢复与固定题复评，不测效果撤销，也未验证任意生成代码不存在外部副作用，因此本章不给出效果撤销成本的实测。
\end{sloppypar}


\subsection{抑制可疑记忆与恢复旧版本}
\label{sec:ch14-functional-forgetting}

恢复历史状态是一种选择，调整当前记忆的选用规则是另一种选择。DREAM预印本提供了后者的具体实例\cite{dream-2610}：给每个外部记忆条目设置可变权重 $w_i$，检索得分由它与查询相似度共同决定。用简化记号写为
\[
 s_i(q)=w_i\,\cos\big(E(q),v_i\big),\qquad 0\le w_i\le1,
\]
其中 $v_i$ 是记忆键的向量。方法对记忆向量施加扰动，以稳定性和相对检索信号形成异常评分，再反复调整权重，目标是抑制可疑记录的影响。在线任务成功信号则按检索得分分配给取回的条目，形成巩固信号；这是一种效用分配规则，并非逐条验证其因果贡献。

该方法不要求物理删除记录，也不更新基础模型参数。它改变的是检索如何选用现有内容，不把状态恢复到某个旧快照。因而评价时应同时检查攻击是否仍能触发、正常任务所需的信息是否保留，以及探查和后续执行的资源消耗。一个较低的攻击成功率不能直接充当文件恢复率、数据擦除率或业务补偿成功率。

作者的威胁模型假设大多数记忆正常，诊断还依赖所用扰动信号能区分可疑记录；其限制讨论承认了解探针的攻击者可能规避诊断\cite{dream-2610}。这个实例的教学价值在于把“修复了记忆”落到明确操作：究竟恢复了内容、改变了访问，还是降低了某种触发机会。采用哪种操作，要由故障位置和相应对照决定。

MemSecBench预印本在其遗忘任务中，把选择性修复定义为两个条件同时成立：目标恶意语义被移除或中和，不再作为有效指导存在，且全部指定正常记忆的语义仍保留（\cite{memsecbench-2607}，原PDF第5页的评价定义及附录E.4）。清空记忆存储即使去掉了恶意内容，也会丢掉这些应保留的信息；仅降低权重，也不保证恶意指导已经失效。原文还指出，最终快照中保留了正常语义，不保证它随后被正确检索或用于任务。对应到本章的恢复设计，目标版本是否还原、应留与应删的内容是否满足要求、恢复后的任务表现如何，需要分别验收。

\subsection{回归题是怎样选出来的}
\label{sec:ch14-selection}

“过去做对的题，现在还会不会做”是合理的工程问题；“更新对原有任务分布的平均影响”则是另一个问题。前者先按过去结果筛选题，后者需要保留事先规定的题目范围。把两者混在一起，容易把随机波动解释成遗忘。尤其当旧版本只运行一次、只有成功题进入回归集时，筛选时的旧分数必然处在顶部，复测却仍可能失败。

考虑一个明确的教学假设：每次无更新运行都以\calcnum{0.5}的概率独立通过同一道题。只保留首次通过的题以后，筛选记录里的旧分数是\calcnum{1}，无更新复测的期望分数仍是\calcnum{0.5}，所以表面下降量的期望也是\calcnum{0.5}。这里没有任何持久更新，却已经出现了正的 $D$。这个算式不描述本书模型的实际成功概率，只说明按成功筛选如何产生向均值回归。

因此，无更新对照不能只是报告表里的一行装饰。应让它在同一批入选题上、与更新臂相同的复测条件下运行，并区分两个量：从筛选记录到无更新复测的下降，以及更新臂相对无更新复测的额外变化。设筛选时旧分数为 $S_b$、无更新复测为 $S_c$、更新后为 $S_u$，则原始下降量为 $S_b-S_u$，相对同期对照的下降量为 $S_c-S_u$。后者减掉的是一个明确定义的对照差，并不自动消除所有时间漂移或题目选择问题。两项都应报告，不能把修正后的量悄悄改名成原始 $D$。

如果研究目标还包括可能的改善，最好另留事先固定、未经成功筛选的能力评价集。否则，单次筛选记录上的满分没有向上空间，更新再好也无法在这个基准上表现为负的 $D$。这不说明改善不可能，而是评价装置限制了可见方向。历史成功题可以继续服务于发布回归检查，固定评价集则承担更广泛的能力变化估计；它们的题目来源、分母与用途应分别保存。

E12 旧草案中“区间包含零即可反对普遍退化”的说法也需要纠正。区间跨零兼容小幅退化、无变化与改善，是证据未定。若要支持损失不超过某个可容忍界，应先确定界值并使用相应的单侧或等效性设计。若要反对严格的普遍退化命题，需要有效证据支持适用范围中的反例。不能在看过结果后把容忍界、更新幅度档位或入选条件改到有利的位置。

\subsection{至少一次成功与多次全部成功}
\label{sec:ch14-repeat-consistency}

回归集确定了检查哪些题，重复运行的指标还要说明怎样算成功。允许尝试多次、找到一次正确答案，与要求多次都可靠完成，是不同目标。$\tau$-bench把后者记作 $\mathrm{pass}^{k}$，把至少一次成功记作 $\mathrm{pass}@k$；其定义针对同一任务的独立同分布重复，再对任务求平均\cite{tau-bench-2406}。原基准固定任务意图与初始数据库，用户模拟器和Agent的消息采样仍会产生不同对话。

设题 $i$ 重复 $n_i$ 次，有 $c_i$ 次通过，且 $1\le k\le n_i$。按原文的组合估计，在 $m$ 道题上可写为
\begin{equation}
 \widehat{\mathrm{pass}^{k}}=\frac1m\sum_{i=1}^{m}\frac{\binom{c_i}{k}}{\binom{n_i}{k}},\qquad
 \widehat{\mathrm{pass}@k}=1-\frac1m\sum_{i=1}^{m}\frac{\binom{n_i-c_i}{k}}{\binom{n_i}{k}}.
 \label{eq:ch14-consistency}
\end{equation}
不足 $k$ 次成功或失败时，对应组合数取零。组合分式是在已有记录中枚举 $k$ 次子集的估计量，不是给每次执行额外增加重试后再挑最好结果。

用一题\calcnum{5}次运行、\calcnum{3}次成功作教学算例。取 $k=\calcnum{2}$，全部成功的两次子集占 $\binom{3}{2}/\binom{5}{2}=\calcnum{0.30}$；至少一次成功则为 $1-\binom{2}{2}/\binom{5}{2}=\calcnum{0.90}$。同一组记录可以同时支持这两个读数，差别在于验收要求。任务间也须先分别计算再聚合，不能把总体平均通过率直接取 $k$ 次幂，代替不同难度任务的平均一致性。

同一个“三次全过”还可能是入选条件。第~\ref{ch07}章的 CompAct-Bench 先取三次运行均成功的题，再从满足长度要求的轨迹中设定续跑边界\cite{compact-bench-2611}。假设题库中两类题各半，每题各次结果独立，两类题的单次通过率分别为 \calcnum{$4/5$} 和 \calcnum{$2/5$}。若只按三次全过筛选，两类题的入选概率之比为 $(4/5)^3:(2/5)^3=8:1$。抽一道入选题，在同一条件下再做一次，通过概率为 \calcnum{$(8\times4/5+2/5)/9=34/45$}。筛选记录里的通过率为一，原题库的平均通过率则是 \calcnum{$3/5$}。这三个数各有自己的题目范围与观察时点。

这些定义还要求复测条件清楚。若一次运行的经验被下一次继承、初始环境没有复位，或模型与服务条件持续变化，就不能直接沿用独立同分布解释。对于自进化系统，可以冻结一个发布版本检查其重复表现，再用明确的更新对照研究能力变化；重复一致性、旧能力退化与更新归因分别回答不同问题。

\subsection{噪声不是一个随手相减的常数}

旧辅助实现把无更新重复间下降量的标准差叫作“噪声底线”。标准差描述散布，不能代表偏移本身。若每一次无更新复测都出现同样的下降，标准差可以为零，而平均下降并不为零。离线构造检查中，两次旧题全通过、两次复测全失败，平均下降为\calcnum{1}，这个所谓底线却为\calcnum{0}。显然不能因为底线为零，就认为所有正下降都来自更新。

也不能用“下降大于一个标准差”代替置信区间或预定发布规则。标准差的样本单位是题、运行还是整批序列，影响其解释；只有一次重复时，把标准差返回零也不能证明随机性为零。工程触发线可以按可接受损失设置，而科学归因需要匹配设计及不确定性估计，两者有不同用途。系统可以因观测到风险而先回退，同时保留“是否由更新造成尚未确定”的结论。

旧 E12 在下一臂读取噪声字段时会报错：该字段要等全部臂结束才写入。附录\ref{app:repro}的入口核对记账摘录的来源，读取当时保存的缺字段观察。修正读写顺序可以让实验继续；“噪声底线”的统计解释仍需独立检验。

\subsection{更新幅度、故障注入与恢复的三个边界}

追加一条、追加多条、再整体重写，可以作为三种操作方案，却不天然构成同一数轴上的纯幅度变化。重写同时可能改变内容、顺序和适用条件；更多条目也可能来自不同来源任务。要将差异主要归到更新幅度，至少应共享同一份已冻结的候选池，让小档是大档的预定子集，并把内容重写作为另一因素。若仍比较三种复合方案，结论就应限定为方案差异，不将相关趋势解释成幅度的普遍因果规律。

文件字节净增量同样不等于修改量。等长替换可能完全改变意义而净增零；增加换行可能增加字节却不改变行为。应同时保存前后工件、条目级增删改以及实际可复核的差异。程序说“重写了全部条目”只说明调用了某个路径，仍须核对输出是否真的变化、是否丢失条目、是否引入新内容。

故障注入还需要验证故障确实发生。把来源任务的通过标签翻转后交给模型，模型可能忽略它、质疑它，或生成一条与真假标签无关的泛化建议。因此“翻了标签”不能自动登记为“已写入一条与事实相反的经验”。注入记录应包含预定错误命题、实际候选、发布工件和受影响接口；未成功植入错误与成功植入但未被检测，是不同结果。本书 E02 的反馈标签扰动也应按这个边界解释。

恢复同时涉及状态与行为。哈希相同确认指定工件恢复；行为复评还依赖外部服务、缓存与实际加载条件。恢复率的分母应明确为更新、恢复操作或复测题；未执行恢复的项，不能仅因没有触发回退线就算恢复成功。E12 v2据此固定共同候选、同期对照及全部更新臂的恢复分母，装置边界见附录~\ref{app:lab}，尚无真实模型恢复率。


\subsection{快照恢复为何需要一套完整事务}
\label{sec:ch14-recovery-transaction}

恢复一份列表在离线测试里可以很简单：读取旧内容，写到新文件，再比较字节。但正在服务的系统还可能有缓存、索引、并发任务、发布指针和外部依赖。若只恢复列表而其他对象继续引用新版本，接下来的任务仍未回到旧状态。因此“可恢复”首先要定义恢复对象与一致性边界。

\subsubsection{把发布对象列成清单}

对于纯文本经验库，最小发布包可能包含条目文件、检索索引、渲染模板与活动版本指针。对于技能库，还可能包括依赖声明和接口契约；对于工作流，还包括执行图与节点实现。哪些对象一起生效，就应考虑哪些对象需要一起恢复。不能只选容易备份的对象，随后把整个系统称为已回滚。

快照还需要描述生成条件。只保存内容摘要值可以检查某个文件是否一致，却不能恢复缺失文件；只保存文件又没有版本与依赖关系，可能无法确定它属于哪次发布。可复查的快照应包含实际内容或可访问的持久对象，以及将它们关联起来的清单。清单本身也应有标识，并在恢复前检查对象是否齐全。

恢复过程可以分成准备、校验与切换。先把目标版本写入新的位置，检查内容、依赖与基本加载是否符合预期，再切换活动版本引用。这样至少避免一半旧文件、一半新文件被读到的明显混合状态。具体原子性取决于存储与运行环境，本书没有验证某个分布式实现，因此这里只说明工程要求，不声称一个文件替换足以覆盖所有服务。

如果正在执行的任务已经加载了新版本，应预先定义是允许它完成、终止后重试，还是把后续步骤切到旧版本。中途切换可能让同一个任务前后使用不同契约；直接重试又可能重复外部动作。最容易审计的记录是每个任务绑定的实际版本，以及恢复时对在途任务采取的动作。没有这两项信息，就难以解释恢复后的第一批成绩。

\subsubsection{恢复失败也必须是一个可见状态}

恢复操作可能在读快照、写目标文件、校验依赖或切换引用时失败。若日志只有“开始恢复”与“恢复成功”，中间崩溃会留下无法解释的空白。应保存每一步已经完成的事实，以及当前对外可见的版本，不能根据一个启动事件就计入成功分子。

每次重试应另存记录，以便判断第一次在哪一步中断。E12 v2 的运行器保留着进程内的基线记忆 $M_0$。恢复时，它把这份记忆写入新文件，核对字节后重新加载，再用加载到的记忆做题。离线检查覆盖了这个过程。若进程已经退出，只能依靠先前留在磁盘上的快照重启，或需要把在线请求切回旧版本，还须检验相应的启动与切换路径。

恢复目标也不能总是“上一个版本”。上一个版本可能已经受污染，只是尚未被发现；更新之间还可能存在依赖，撤掉某条经验却留下依赖它的工作流，仍会出错。可以依据已知通过范围选择一个检查点，再核对之后的派生对象。选择检查点的依据应记录为证据与决策，而不是在事后把得到最好复测分数的版本倒称为原定恢复目标。

恢复后的验证有两层。状态层核对预定工件的内容，行为层逐题核对复评结果。若写好了恢复文件，复评却因预算耗尽而中断，文件字节仍可核对，未完成的题只能记为缺失。附录~\ref{app:lab} 的“先保存派发记录，再允许执行”一节给出了离线示例。第二道测试只写入了派发预留，重开日志时没有判决。后来按演示设定完成清理对账，这道题仍记为缺失。

\subsection{恢复率的分母决定了它的含义}
\label{sec:ch14-recovery-denominator}

一个“恢复率”至少可能指文件操作成功率、受损更新的恢复覆盖率、复测任务通过率或达到某个质量界的版本比例。这些量没有共同分母。若报告者只挑选执行过恢复的成功案例，会遗漏没有发现的有害更新、被预算挡住的恢复，以及恢复过程中的基础设施缺失。

可以先为每次更新保存一条生命周期记录：是否发布、是否触发工程回退线、是否尝试恢复、状态层是否恢复、行为复评是否完整。每个阶段都有自己的分母。没有触发的更新不属于“触发后恢复成功率”的分子；未执行恢复也不能因为任务暂时正常就算恢复成功。缺失复评应单列，而不是随意并入成功或失败。

E12 v2 冻结设计对所有更新臂都安排一次恢复与复评，正是为了不只观察跌得最重的臂。这样增加的成本属于恢复评价，必须预先计入预算。实际部署可以只对达到业务线的更新恢复，但那时报告的结果是一个被触发规则筛选过的子集，不能代表所有更新。实验设计与生产策略可以不同，报告需要分别说明。

E12 v2 把整批固定题作为判断单位。冻结的小试用四道题，每题判一次通过或失败。基线通过率减去更新后通过率，是原始下降量 $D_{\mathrm{raw}}$。工程回退线定为 $D_{\mathrm{raw}}\geq 0.10$。四题中净少通过一题就是下降 $1/4$，已经触线。运行器在整批评价结束后才判断是否触线，更新臂的完整记录还要等恢复后复评结束才写入。凭逐题结果可指出哪道题转为失败，却无法据此给出逐题报警时刻。工程上的容忍度与统计不确定性须分别说明。

用一个算例看总分会掩盖什么。假设基线只通过甲、乙，更新后只通过乙，恢复后通过乙、丙。复评总分回到原来的 $2/4$，更新时失去的甲依旧失败，恰好由新通过的丙补足。若问题是“原先失去的题找回了多少”，分母应取基线通过、更新后失败的题，分子取其中复评又通过的题。本例为 $0/1$。若复评通过甲、乙，则同样的总分对应 $1/1$。若更新时一道旧题都没丢，这一比例没有分母。每题一次的判决只能描述这一次变化。

对于已经发生的外部效果，还应另列补偿范围。撤回错误记忆可以阻止它继续影响新任务，却不会自动取消已经发送的请求或修复已经写入的外部记录。补偿是一个新的动作，可能有自己的权限、费用和失败风险；原动作与补偿记录都应保留。E12 的文本工件实验没有评价这部分，不能从文件恢复率推导实际业务恢复率。

\subsection{影子运行与分阶段发布的不同证据}
\label{sec:ch14-shadow-evidence}

回归题固定了可重复检查的范围，真实服务却不断出现新输入。影子运行可以让候选处理一份输入副本，同时继续由当前版本提供正式输出；分阶段发布则让候选实际承担一部分工作。两者都能补充回归题，但暴露的信息与风险不同，不能只在图中画成发布前后的两个勾选框。

\subsubsection{影子输出不生效，成本仍然发生}

在理想的影子模式中，候选可以生成答案或动作计划，但这些动作不会对正式环境产生效果。若候选调用了会写数据库、发送消息或修改远程资源的工具，它就不再只是旁观执行。需要为影子任务提供只读接口、可控的替代服务或隔离副本，并记录哪些动作被模拟、哪些真实执行。本书没有部署这样的生产系统，所以这是一组实现要求，不是已完成的可靠性结果。

影子运行的模型调用、检索、测试和存储照常产生资源消耗。它的输出不进入用户服务，不表示其费用可以从更新成本中删去。若同时运行多个候选，资源争用还可能拖慢正式路径；这种影响应与候选自身时延分开观察。仅报告最终被选中候选的推理费用，会漏掉影子筛选阶段的整体代价。

影子比较也不总能复制正式环境的反馈。正式版本的某个动作可能改变后续输入，而影子候选由于没有实施自己的动作，只能继续看到由正式版本造成的环境轨迹。这类似第\ref{ch13}章的静态回放限制：能够比较当下输出，并不等于观察了候选独立运行后会面对的全部未来。对于多步交互，应明确哪些后续状态是共享的、哪些是候选专有的。

把正式答案与影子答案交给评委，也会引入评委误差。若评委看到版本名、解释长度或展示次序，评分可能受这些因素影响；若输出涉及可执行任务，应尽量使用与目标一致的外部检查。无法获得这样的检查时，可以保留盲化和配对记录，但仍须承认判断误差，不能因为流程叫“影子评测”就自动提升证据等级。

\subsubsection{少量放行不自动代表风险小}

分阶段发布通常从一部分工作开始，再根据质量与资源表现增加范围。但抽到哪些任务，比单纯放行比例同样重要。如果初期只处理简单、短小或低价值任务，表现良好并不证明复杂任务也安全；若按用户群或项目分配，组间差异也可能影响结果。需要说明分配规则、覆盖范围与未覆盖部分，而不是只报一个流量比例。

对于有持久记忆的系统，版本之间还可能互相污染。当前版本与候选若写入同一个经验库，候选生成的内容可能被当前版本取走，当前版本就不再是未更新对照。共享索引、缓存与后台整理也可能造成类似混合。实验可以允许某些共享对象，但必须明确定义哪些状态属于每个版本，哪些是公共环境；否则“候选与当前版本的差异”缺少清楚对象。

退出条件应在放行之前确定。它可以包含明显的接口错误、预定损失容忍线、资源异常或缺失率，但不能在看到不利结果时临时延长观察直到回升。若需要重复查看流式结果，应为这种查看规则设计相应评价，而不把每个时点的普通区间都当成一次独立终检。当前书中小型实验没有实现在线持续检验，正文不借用实时监测的名字宣称已解决这一问题。

候选退回后，已收集的事件仍应保存。那些事件可能暴露了过去回归集没有覆盖的条件，可用于未来开发，但不能立即并入当前最终测试后重新宣布候选通过。把事故题加入开发回归集是合理维护；是否还保留一个未用于调试的评价范围，是另一个需要安排的步骤。两者混在一起，会让测试集逐渐变成已经记住的事故目录。

\section{长期退化的三种漂移}
\label{sec:ch14-drift}

\begin{sloppypar}
双轨只有在持久保存、访问保护和备份实际成立时，才有助于防止原始记录丢失。它仍留下另一类问题：
记录都在，系统的行为却慢慢偏离了。一份记忆治理的工作把这类偏离分成三种\cite{ssgm-2603}，
这些变化可以逐步累积，也可能由一次重要的改写突然显现：
\end{sloppypar}

\begin{itemize}
  \item \textbf{语义漂移}。反复的有损压缩（原文 $\rightarrow$ 摘要 $\rightarrow$ 更高层反思）
        逐步剥掉细节。该文给的例子是一条原本温和的偏好被逐次重写成强偏好，
        最终导致一次明确的偏好违背。
  \item \textbf{过程漂移}。次优的执行路径被强化并固化进过程记忆，
        之后即使有更好的做法也进不来。第 5 章讨论的求解器特化可能与它发生联系，但特化本身不等于过程退化。
  \item \textbf{目标漂移}。记忆中累积的偏差让行为慢慢偏离原始指令。
\end{itemize}

\begin{sloppypar}
该文把语义漂移写成可测量的形式：以当前记忆与理想语义目标在嵌入空间中的余弦相似度定义漂移
$\delta(M_T, K_{\text{true}})=1-\mathrm{sim}\big(E(M_T), E(K_{\text{true}})\big)$，
原文进一步区分理想目标与可获得的原始账。情景账提供校正材料，不能自动充当无误真值；嵌入相似度也只是一个表征层的代理指标，须与具体事实、范围和行为检查配合。
\end{sloppypar}

% ---------- 7. 闭环退化 ----------
\subsection{SSGM 的漂移界能支持哪一种结论}
\label{sec:ch14-drift-bound}

SSGM 给出了一个有条件的漂移界：每次合法摘要或整合引入的语义误差不超过某个步长界，并且对账算子能够把剩余误差压回一个不随总时长增长的范围；若每隔固定步数执行对账，误差累积的主要项就与对账间隔有关，而不是一直随总运行长度增长\cite{ssgm-2603}。理解这个结论，关键在于保留它的假设。

用本书的简化记号，设一次对账后相对固定语义目标的期望误差不超过 $C$，之后每步使这一误差的期望增长至多 $\epsilon$。若两次对账相隔至多 $N$ 步，对增量逐项求和才得到
\[
 \mathbb E[\delta_t]\le C+N\epsilon.
\]
这是对原文 Theorem 1 的条件性解释：原文讨论期望漂移，并在证明中采用可加误差上界。若另有逐路径的同类增量界，才能给出相应的逐路径结论。这里没有证明任何具体摘要器、对账器满足这些前提；对账漏掉事实或引入新错误时，残余界也需要重新验证。

\subsubsection{相邻摘要都接近，首尾能相差多少}

原文式 (4) 使用 $d(u,v)=1-\cos(u,v)$。这个量有一个容易忽视的性质：对非零向量，它不是满足三角不等式的距离。取二维单位向量，方向依次为\calcnum{$0^\circ$}、\calcnum{$60^\circ$}、\calcnum{$120^\circ$}，记为 $u,v,w$，则
\[
 d(u,v)=d(v,w)=\calcnum{\tfrac12},\qquad
 d(u,w)=\calcnum{\tfrac32}>d(u,v)+d(v,w).
\]
若每次只量相邻摘要，仍然可以给出一个几何上界。设 $u_0,\ldots,u_k$ 都是同一实内积空间中的非零嵌入向量，且每一步 $d(u_{i-1},u_i)\leq\epsilon$，其中 $0\leq\epsilon\leq 2$。将向量归一化后，相邻方向的夹角至多为 $\theta=\arccos(1-\epsilon)$。单位球面上，两点之间的最短弧长等于它们的夹角，夹角满足三角不等式\cite{schubert-cosine-2021}。把这条不等式逐步相加，首尾夹角至多为 $k\theta$；任意两个方向的夹角又至多为 $\pi$。余弦在这段区间上单调递减，因而
\[
 d(u_0,u_k)\leq 1-\cos\!\left(\min\{\pi,k\theta\}\right).
\]
上例每步最多转 \calcnum{$60^\circ$}。连续两步朝同一方向转动，首尾的余弦差就能达到 $\calcnum{\tfrac32}$；允许三步时，终点已经可以指向反方向，余弦差可达 $\calcnum{2}$。这个界描述首尾嵌入相隔多远。要将它解释为相对理想语义目标的误差，还需知道起点与该目标之间的关系。零向量没有可供比较的方向。

原始账与真实目标之间也可能有差距。账本可以忠实记录用户当时说了某件事，却不说明事情本身为真；不同来源还可能互相矛盾。只追加机制保证旧记录不被静默覆盖，不能保证记载内容正确。因此对账中的“回到来源”可以修复摘要失真，却不能自动修复来源误报。这是来源完整性与事实正确性的区别。

误差度量同样重要。嵌入相似度可能对整体主题很敏感，却对否定词、日期、数量或权限范围不敏感。一份摘要把“可以”改成“不可以”，主题向量仍可能很接近，而行为后果已经反转。一个表征距离保持较小，不等于所有关键事实与决策风险都有相同界。工程上应在总体相似度之外加入关键字段与行为测试，并说明每项检查覆盖什么。

对账频率还有成本约束。更短的间隔可以缩短上述假设下的累计窗口，但会增加扫描、重建和验证工作。若对账速度赶不上更新速度，待处理队列还可能增长。不能只在证明里让 $N$ 尽量小，就宣称同时获得稳定与低成本；实际系统要测完成延迟、队列长度及任务阻塞情况。SSGM §6.3 将延迟与一致性的取舍列为后续研究假设，而不是给出了普遍免费治理的实测证明。

这也解释了为何本书把双轨设计作为待检验机制。它提供可追溯与可重建的条件，有助于提出清楚的实验；其效果仍取决于来源、校验器、对账算子及恢复实现。E12 的小型文本快照检查覆盖的是其中一个很窄的切面，不能替整套治理框架背书。

\section{闭环退化的三类机制}
\label{sec:ch14-collapse}

\begin{sloppypar}
闭环还可能引入另一组风险。一份 RSI 综述列出四类失效模式，下面先讨论其中前三类\cite{rsi-survey-2607}，各自留有不同的签名：
\end{sloppypar}

\begin{table}[htbp]
\centering
\caption{三类闭环退化：机制、可测签名、本书的干预点}
\label{tab:ch14-collapse}
\begin{small}
\begin{tabular}{>{\raggedright\arraybackslash}p{2.4cm}>{\raggedright\arraybackslash}p{4.4cm}>{\raggedright\arraybackslash}p{3.4cm}>{\raggedright\arraybackslash}p{3.2cm}}
\toprule
类型 & 机制 & 可测签名 & 本书的干预点 \\
\midrule
自证循环 & 共享机制可能造成错误相关 & 错误重合与置信度校准；独立证据覆盖 & 反馈源分级；接地优先 \\
\midrule
塌缩 & 部分递归自训练设置中分布尾部丢失 & 输出分布熵；与真实数据混合的比例 & 接地频率 \\
\midrule
多样性塌缩 & 共进化中\textbf{任务分布}收窄 & 任务来源与类别覆盖；固定题分数另报 & 任务族轮换；新颖性预算 \\
\bottomrule
\end{tabular}
\end{small}
\end{table}

\begin{sloppypar}
任务分布收窄与固定任务上的分数先升后降，是不同现象。所引综述在同一讨论中列举了这两类研究，但前者不能仅由后者推断；固定任务上的下降也可能有其他机制\cite{rsi-survey-2607}。因此应结合任务来源、类别覆盖、生成模板与长期表现，而不是只画总成功率。本书尚未为 E11 的有限固定序列建立一个可裁判的任务分布熵指标，不把“要求报告”冒充已测结果。一个高熵标签直方图也不能证明问题语义多样。
\end{sloppypar}

\begin{sloppypar}
还有第四类更难：\textbf{框架锁定}。
前三类可为循环内部诊断提供线索——置信度耦合的奖励、正在塌缩的输出分布、
正在收窄的课程，都在循环已经追踪的量上留下签名；
框架锁定可能不被原有指标覆盖。综述提出，一个循环可以在现有指标上表现为接地、诚实、多样与稳定，
却仍在优化一个\textbf{早就不值得优化的目标}（原文称之为"well-posed wrong question"）。
检出它需要一个位于循环自身判据之外的视角，而这正是闭环所缺少的东西。
本书不声称解决了它\cite{rsi-survey-2607}，只在第 16 章把它列为开放性边界。
\end{sloppypar}

% ---------- 8. 最小接地 ----------
\subsection{最小接地兑换率}
\label{sec:ch14-grounding}

\begin{sloppypar}
塌缩那一类留下的是本学科目前最实际的一个开放问题。
该综述综合所引研究讨论了纯闭环退化的风险，并指出实用系统可以引入外部信息\cite{rsi-survey-2607}。这不是本书证明了任意闭环必然退化。相应的开放问题是
“最少需要多少外部接地就够”，而\textbf{本书已核对材料尚未确立一个跨任务通用的兑换率}。
\end{sloppypar}

\begin{sloppypar}
本书为它准备了候选实验 E02b（候选命题 P8）：
在接地频率的若干档位上测收益曲线，看是否存在远低于“每轮全接地”的最小接地量。
它属于主线之外的扩展研究，当前尚未启用。可在协议中先检验该曲线是否具有可区分的形状，正文不提供未发生的兑换率实测。
\end{sloppypar}

% ---------- 9. 条件变化 ----------
\section{让 E12 的新比较可以被复查}
\label{sec:ch14-e12-v2}

E12 v2将固定评价、共同候选池与恢复安排实现为独立组件，离线反例和旧代码身份单独保存。下文说明这些条件如何形成可复查的新比较；对应真实模型试验尚未运行。

新的固定评价在模型运行前规定题集，不按旧版本是否通过筛题，从设计上保留观察损失与改善的空间。各臂必须具有完全相同的任务标识集合；缺题、重复题或非布尔的基础设施状态会被拒绝进入完整比较。旧版共同基线分数为 $B$，同期无更新对照为 $C$，更新臂为 $A$，则分别报告
\[
 D_{\mathrm{raw}}=B-A,\quad D_{\mathrm{control}}=B-C,\quad
 D_{\mathrm{adjusted}}=C-A.
\]
最后一项是前两项之差，表示相对同期无更新对照的描述性变化。它不是从标准差中扣出来的“去噪分数”，也不能自动消除执行顺序或服务漂移。

更新内容从一份固定候选池取预定前缀：追加一条与追加三条共享第一条；重写比较则在追加三条的基础上替换指定旧条目，单独作为因素。这样避免每臂独立生成完全不同来源经验后仍称只改变幅度，但也没有消除内容本身的影响。结论最多限定于该候选池与该追加路径，不能从三条经验推出全部记忆系统的单调退化规律。

构造性污染采用明确固定的错误断言，并在运行前保存可执行反例。例如本书组件使用“Python 的 range(a,b) 包含右端点 b”作为错误事实，用具体整数输入检查右端点实际未出现。该检查能证明预定语句在所指接口下为假，不能证明模型自然会产生它，也不能证明写入后模型一定采纳。后续仍需核查发布工件与实际上下文。

幅度检查分开保存新增、删除、文本修改、来源元数据修改与顺序变化。字节对齐给出另一组可重复计算的差异量，但它不是最小编辑距离，更不是语义距离。格式变化可能产生很多字节差，极短的否定词修改也可能造成重要行为变化。因此幅度报告服务于解释操作，不替代旧能力评价。

E12 v2 把共享准备费用记一次，共同基线评价另列。每个更新臂的专属小计合并了更新后做题和恢复后复评的用量。要回答“复评又花了多少”，需按调用日志中的回合编号挑出该臂的复评请求，核对已知用量，并单列失败传输中用量未知的部分。各臂小计已经包含复评，合计全程用量时不可再加一次。

预检记录保存了来源题、评价题、执行顺序与请求上限。冻结的小试计划只安排一次重复，并规定无更新对照先执行；这限制了时段控制与统计精度。清单约束的是保存时的源码版本，共享执行器更新后，应在执行前重新核对文件身份与依赖绑定，不能继承旧预检的通过状态。对应真实模型试验尚未运行。P6分别研究哪些更新引起退化、幅度怎样相关，以及怎样恢复；这三项需要不同的估计量与对照。


\section{一次事故怎样改变后续更新规则}
\label{sec:ch14-incident-learning}

“恢复以后不要再犯”需要比增加一句警告更具体的工作。事故可能来自错误来源、错误概括、错误检索、权限越界或验收遗漏，不同原因对应不同修复位置。若每次都只往记忆里追加“下次更小心”，规则会越来越长，真正的缺陷却仍留在执行路径上。

\subsection{先定位最早偏离，再选择修改对象}

可以沿事件链依次检查：原始输入是否符合预期；候选是否忠实保留条件；准入是否按规则执行；发布对象是否完整；后续任务是否读取正确版本；验收结果是否反映目标。找到最早有充分证据支持的偏离点，有助于避免把后果当成原因。没有完整记录时，应把原因写成待验证假设，不能凭模型给出的流畅解释直接定案。

Liu审计HAL评测平台公开的TAU-bench Airline轨迹时发现，通用代理Generalist的对话结束方式常未触发评分，未评分回合被归档为零分。评分器为构建参照状态而回放的标准动作又混入轨迹日志。作者区分动作来源，并用固定基准版本对Agent自身的动作离线复评。修正后的汇总保留已评分回合的原结果，为未评分回合补入复评结果\cite{silent-instrument-2611}。

这次漏评把多少成功算成了失败？在标准答案含至少一个非终止动作的题目上，Generalist 共留下 \cited{613} 个回合。一道题可以对应多个回合。按作者的勘误逐回合表复算，原有 \calcnum{429} 个零分，其中 \calcnum{130} 个是评分器判定失败，另有 \calcnum{299} 个根本没有评分。后者补判出 \calcnum{77} 个成功。若问漏评组有多少通过，分母是 \calcnum{299}；若问总体成功率改变多少，分母仍是全部 \calcnum{613} 个回合。保留原先判定成功的 \calcnum{184} 个回合，总体成功率便从 \calcnum{$184/613$} 升至 \calcnum{$261/613$}，增加约 \calcnum{12.56} 个百分点。改变的是计分记录，代理完成的动作仍然是原来那些。

例如一条正确经验被错误地分配给另一任务族，修改经验措辞未必能解决问题，检索过滤与适用条件才是优先检查对象；若候选没有满足事实要求却通过格式门，则需要补充准入机制或缩小发布范围；若工具反馈被解析错，先修解析器通常比训练模型忽略某个固定字符串更直接。这里的优先级是工程诊断，不是未经对照就证明哪种修改最有效。

修复以后，应保留一个能够复现原缺陷的局部检查，以及一个与目标行为相关的完整任务检查。前者证明已知路径受到约束，后者检查修复是否破坏正常工作。单纯为新实现写一个返回相同字符串的测试，没有验证原来的失败条件；只在原事故题上通过，也不能证明其他条件下没有回归。局部与端到端证据应互相补充。

E11 的共享预算修复就是这种思路：不仅检查计数器会增加，还让脚本后端在真实本地任务中用尽配额，确认未完成序列不会被填成失败结果、候选构建也不能绕过总账。E12 的等长替换检查则验证净字节差为零时仍记录文本修改。这些离线检查针对测量装置的实际风险，而非模型能力；它们进入来源索引，避免后来把“测试通过”误写为命题得到支持。

\subsection{保留失败证据，不把修复写成成功历史}

一次实现修订可能改变模型提示、候选准入、预算寿命或统计单位。只要这些变化影响比较对象，就需要明确的新版本，而不能悄悄替换旧代码后继续使用旧 run ID。旧轨迹仍有价值，可以用于发现缺陷、重算可识别的账目或展示工程反例；但它不能被自动升级为满足新协议的结果。

当错误只在报告层，例如把同一项费用重复累加，原始调用记录可能足以重建新的分析结果。此时可以保留原始数据、另出分析版本，并说明改了什么。若错误发生在执行层，例如某臂实际多拿到信息、某次失败后预算被重置，事后重算并不能让当时的实验满足新的边界。需要区分“可修复的分析”和“必须另做的行为比较”。

共享组件的修订也会影响后续实验的可比性。如果一个新的解析器改善了多行动回复的处理，它可能同时改变成功率与消耗轮数；将旧基线与新方法直接比较，差异就包含了解析器变化。为避免这一点，应把执行器版本作为设计条件冻结，或者在新版本下重新安排明确获准的配对比较。不能把这些差异全部归入持久学习方法。

事故后的规则还需要生命周期。某个绕过措施可能只针对旧工具版本，若永久保留，会阻碍以后正常行为。可以为这类规则记录起因、适用条件、验证范围与何时重新检查。退役也需要证据：删除前确认相关故障条件是否仍存在，删除后观察必要的回归范围。这样既不把每次事故都变成永久补丁，也不因规则显得陈旧就随意丢弃。

最后，应把证据状态传递到书面结论。已经定位的实现缺陷可以明确写出；尚未复现的风险仍是可能性；缺失的成本记未知；未运行的能力比较保持未定。可靠性写作不能以更强的语气补偿证据不足。对一个自进化系统而言，能够准确说出哪一层已经恢复、哪一层尚未验证，本身就是继续改进所需的基础信息。

% ---------- 6. 漂移 ----------
\section{条件变化：哪些防线在什么时候失效}
\label{sec:ch14-conditions}

\begin{itemize}
  \item \textbf{持久物不便直接解释时}（例如第 6 章讨论的潜记忆），
        文本摘要式审计难以直接使用，但仍可保存状态快照、输入来源和版本记录；若存在可比较的状态表示，也可设计距离指标。
        这些检查需要与端到端行为评价结合，不能仅因状态不可读就断言全部监测失效。
  \item \textbf{反馈只有一个来源时}，自证循环无法用“换一个反馈源”来破，
        可考虑引入独立依据、延迟发布或缩小更新范围；E02b 是探索接地频率的一个候选方向，并非已经验证的唯一办法。
  \item \textbf{任务分布由系统自己产生时}，可能出现多样性塌缩或框架锁定，
        而任务分布熵这个签名本身也被系统控制——需要检查类别与生成过程，并可用独立题源或固定锚点作对照，不能仅凭自报的熵认定覆盖充分。
\end{itemize}

% ---------- 10. 能说与不能说 ----------
\label{sec:ch14-claims}

这套检查分别回答三个问题：输入是否受到扰动、发布的状态是否改变，以及改变是否影响后续能力。E02的两条标签翻转只直接确认第一项；负向点估计与跨零区间描述各自训练、构建后的方案比较，尚未隔离标签翻转的净效应。回归评价与恢复试验则需要预定题目、同期对照和明确的恢复对象，不能让文件恢复成功代替旧能力复评。

若沿用本章两项噪声比较的任务级差值标准差，按三个百分点半宽作近似精度规划，两者均需\calcnum{490}道配对题，而当前各只有\measured{16}道。这个数来自噪声臂分别减无记忆基线、未翻转测试反馈臂的差值；它不使用第\ref{ch04}章另外两项反馈比较的方差，也不构成完整功效设计。E12和E02b尚无真实结果，双轨、恢复与接地频率的价值仍需在各自明确的条件下检验。

\begin{principlebox}
可靠性机制既应降低错误发生和传播的机会，也应\textbf{让事故可被定位与恢复}。日志与断言只能覆盖它们实际观察和约束的路径。
一条记不清来源的记忆、一次无法定位的回退、一个无法判断是哪一侧被改动的提升——
这三者是同一类失败：观测被切断了。
所以本章的每一条机制都配一个落点：查询账与执行权限分别检查，
双轨落在受保护的原始记录与可核对的快照上，漂移同时检查语义参照、关键事实与后续行为。
\end{principlebox}

\falsifier{P6 的普遍退化主张可被适用范围内可靠的反向证据挑战：若有效配对设计支持某次更新使旧能力改善，即下降量 $D$ 的不确定性范围支持负值，则与“任何更新均退化”不相容。区间含零只表示未定，不能自动证伪；区间下界为正反而支持该范围内的退化。单调性必须另作检验。E12 尚未运行，本章不判决。}

\openquestion{本书尚无 E12 的真实回滚成功率与更新退化实测。事故处理属于完整生命周期成本，不能因为未测就从定义中删除；应保留该项为未知，并报告已有账目的覆盖范围。新版方案规定了可恢复状态、回归题、故障注入及复测条件；执行前还须把当前源码与依赖重新绑定到清单。随后在明确授权的规模内取得真实记录，再判断是否需要修订设计。}

\section*{习题}
\begin{enumerate}[label=14-\arabic*.,ref=14-\arabic*]
  \item\label{ex:ch14-1} 章首的错误标签没有特殊格式。请区分“标签被翻转”“模型察觉矛盾”“错误内容进入记忆”和“后续任务采纳错误”四个主张，为每个主张指定需要检查的记录，并给出一条利用原始执行输出的检测。
  \item\label{ex:ch14-2} 双轨设计若允许受控读取原始记录，需要保存哪些访问与来源信息？请比较“原始内容真实”和“原始内容未被改写”两个主张，并说明摘要重算如何依赖它们。
  \item\label{ex:ch14-3} 表~\ref{tab:ch14-collapse} 的第三类（多样性塌缩）里，
        区分任务分布收窄与固定任务分数下降。请设计两种不同对照，说明为什么单个总分或类别熵未必足以识别原因。
  \item\label{ex:ch14-4} 框架锁定可能不被现有评分规则捕捉。请说明仅提高同一评分规则的执行精度为何未必能发现目标本身不合适，
        以及你会把哪一类外部视角加进系统。
  \item\label{ex:ch14-5} E02 噪声臂的两个区间都跨 0，因此本章不说“污染一定有害”。
        请提出一个范围明确、可被反向证据挑战的命题，分别列出精度设计与功效设计所需的假设。说明为何仅凭现有方差无法得出适用于所有目标效应的样本量，并比较式~\ref{eq:ch14-consistency}的两种成功要求。
\end{enumerate}

\section*{常见误解}
\begin{itemize}
  \item \textbf{“加了隔离就安全了”。}
        权限隔离能限制部分改动路径，仍不能证明所有可写内容都正确；还需验证实际访问范围和候选行为。
  \item \textbf{“污染一定会让分数掉下来，所以看分数就够了”。}
        E02两项噪声比较的区间都跨零；分数未下降不能排除污染，区间跨零也不能确认错误内容已写入或被采用。
        本次比较尚未确定总体差异方向，也未单独隔离各种污染路径的作用。
  \item \textbf{“原始记录留着就行，摘要写坏可以以后再说”。}
        没有原文会限制来源重建；若有完整历史快照，仍可能恢复旧状态。
        回滚的直接对象是可变状态；原始账用于追溯或重建，历史快照用于核对目标版本。
  \item \textbf{“接地越多越好”。}
        外部反馈也可能错误、过期或昂贵；增加频率并不自动保证质量单调提高。本书尚无足够证据确定跨任务通用的最低频率或收益曲线。
\end{itemize}

\section*{进一步阅读}
\begin{itemize}
  \item \cite{rsi-survey-2607}：§5.3 的四类失效模式是本章 §\ref{sec:ch14-collapse}
        整节的骨架，包括“奖励破解只是自证循环的显式特例”与外部接地需求的开放问题；这些是综述归纳，不能代替本书对所有机制的直接复现。
  \item \cite{silent-instrument-2611}，§3.2—3.3分析评分触发与动作来源，公开勘误补足中途插入的标准动作识别。§3.4另讨论评价器版本变化。
  \item \cite{tau-bench-2406}：§3中$\mathrm{pass}^{k}$与$\mathrm{pass}@k$的定义、组合估计及独立同分布前提；第\ref{ch15}章再比较结果状态与过程约束。
  \item \cite{dream-2610}：§3—4的检索权重调节、附录A的适应性攻击边界与附录G的在线巩固信号；读时区分功能抑制、快照恢复与物理删除。
  \item \cite{memsecbench-2607}：原PDF第5页及附录E.4的恶意语义移除与正常记忆保留，区分选择性修复、快照恢复和后续任务使用。
  \item \cite{ssgm-2603}：§6.2 的双基底与 Figure 4，§6.1 的写前校验门；
        “漂移不是一件事”那一节与式 (4) 是本章漂移定义的出处。
  \item \cite{coding-survey-2608}：§5.2 把组件退役与回滚成功率列为应当报告的度量，
        本书把这种建议落实为待执行的记录规范，不能据此宣布全部文献均无恢复证据。
  \item 附录~\ref{app:lab}：隔离边界在实验台上落在哪些字段，
        \texttt{persistent\_writes} 与 \texttt{QueryBudget} 两处断言的实现位置。
  \item 附录~\ref{app:stats}：区间精度、检验功效、抽样范围与证据状态的区别。
\end{itemize}
